\section{Erd\H{o}s Problem \#319: independent continuation}
\noindent Date: 23 September 2026. Source versions read in full: \texttt{319.tex} in \texttt{gpt\_pro\_5.4}.

\subsection*{1) OBJECTIVE AND SOURCE AUDIT}
The problem maximizes the size of a minimal signed vanishing reciprocal
sum, with no proper nonempty zero-sum subcollection. The source settles
the asymptotic density of the singleton-sign construction. Its proof
applies a fixed-$r$ asymptotic with $r=1/p$, although $p$ can depend on $N$.
The following finite inequality fixes that uniformity issue and extends
the obstruction to every bounded size of the smaller sign class.

\subsection*{2) WORK: BOUNDED SIGN CLASSES CANNOT APPROACH DENSITY ONE}
Write $A=P\sqcup Q$ and suppose the two reciprocal sums agree. Put
$h=\min(|P|,|Q|)$ and, by exchanging signs, assume $|P|=h$, $|Q|=m$.
Since the largest possible reciprocal mass of $h$ distinct positive
integers is $H_h$, while the smallest mass of $m$ members of $[1,N]$
is a terminal harmonic sum,
\[
 H_h\ge\sum_{p\in P}1/p=\sum_{q\in Q}1/q
 \ge H_N-H_{N-m}
 \ge\log\frac{N+1}{N-m+1}.
\]
The last inequality follows by integrating $1/t$ over $[N-m+1,N+1]$.
Hence the uniform finite bound is
\[
 |A|=h+m\le h+(N+1)(1-e^{-H_h}). \tag{1}
\]
For $h=1$ this recovers the source's upper density $1-e^{-1}$ without
any parameter-dependent error term. For each fixed $h$, it gives the
upper density $1-e^{-H_h}<1$. Thus a sequence of admissible sets with
$|A|/N\to1$ must have both sign-class sizes tending to infinity.

More quantitatively, using $H_h\le1+\log h$ in (1),
\[
 N-|A|\ge\frac{N+1}{eh}-h-1.
\]
This is nontrivial, for example, when $h=o(\sqrt N)$. It isolates a
necessary growth requirement on the number of positive and negative terms
for any construction improving strongly on the singleton mechanism.

\subsection*{3) ADVERSARIAL VERIFICATION}
The inequality needs only equality of the two masses, so it applies in
particular to minimal signed zero sums. It does not prove minimality
for a proposed new construction. Combining two existing signed relations
would generally violate minimality by retaining a proper vanishing subset.
No matching construction for the fixed-$h$ bound is asserted.

\subsection*{7) FINAL AND REFERENCE}
\textbf{UNRESOLVED.} A uniform finite bound and a larger family of structural
obstructions are proved, but the maximal size of a minimal relation remains
undetermined. The minimality condition was checked in the indexed primary
statement: \url{https://www.erdosproblems.com/319}.

\subsection*{8) COMPLETION ESTIMATE}
\textbf{COMPLETION: 10\%}
This is a subjective estimate of progress toward the entire stated problem, not a probability that the conjecture or an unverified proof is correct.
